The pairing is real, symmetric and nondegenerate on Hermitian matrices, since the real trace is cyclic and $\operatorname{Re}\operatorname{tr}(A^2)$ is the sum of the squared entry norms. In particular $A:(vv^*)=v^*Av$. A Hermitian quaternionic matrix has an orthonormal eigenbasis: maximize its quadratic form on the unit sphere and apply the first-variation condition to obtain an eigenvector, then repeat in the orthogonal complement, which is invariant by Hermitian symmetry. Newton's identities express its determinant as a polynomial of the power sums $\operatorname{Re}\operatorname{tr}(A^j)$ for $1\le j\le m$, giving the determinant and cofactor degrees. This is the Moore determinant convention of \cite{r-2}, Section~1.1.

In the real representation on $\mathbb R^{4m}$ every eigenvalue occurs four times. For $A>0$, $\log\det A$ is one quarter of the real log determinant. Its derivative in a direction $B$ is $\operatorname{Re}\operatorname{tr}(A^{-1}B)$ and its second derivative along a line is $-\operatorname{Re}\operatorname{tr}((A^{-1/2}BA^{-1/2})^2)\le0$. The first derivative is nonnegative for $B\ge0$. These identities yield the positive definite cofactor formula, monotonicity and concavity. In a diagonal basis the cofactor eigenvalues are the products of the other $m-1$ eigenvalues, first for positive $A$, then for all diagonal $A$ by polynomiality. Unitary invariance gives the assertions for $A\ge0$ and for its image in the singular case.

At a vector $(b,w)$ the form of $M-\sigma v_0v_0^*$ is $(w+D^{-1}cb)^*D(w+D^{-1}cb)$ by completing the square. If $\sigma>0$, the real Schur determinant formula in the real representation gives $(\det M)^4=(\det D)^4\sigma^4$, and positive definiteness gives the positive fourth root. Both sides of $\det M=(\det D)\sigma$ are polynomial in $\ell$, extending the result to all $\sigma$. To prove the rank-one formula, extend the unit vector $v$ to a quaternionic orthonormal basis and use its unitary matrix $U$, so $U^*v=e_1$. Apply Proposition~\ref{ext:alesker-determinant} to $U^*AU$ with $T=\operatorname{diag}(t,0,\ldots,0)$. The only terms of the diagonal expansion are the unchanged determinant and $t$ times the principal minor on $\{2,\ldots,m\}$, which is $\det(A|_{v^\perp})$. This proves the perturbation identity for every Hermitian $A$. Differentiating at $t=0$ and using the real trace pairing yields $v^* C(A)v=\det(A|_{v^\perp})$.
